\subsection{Wasserstein gradient flow on product measures}
\label{app:product-measure-wgf}

Before analyzing the matrix-decomposition functional, we record a simple form of
Wasserstein gradient flow restricted to product measures.  This is useful when the
optimization variable is constrained to remain factorized across parameter blocks,
so that admissible perturbations move each marginal separately.
Related Wasserstein product-space and multi-distribution gradient-flow
structures are studied by
\citet{chen2016wassersteinProduct,conger2025coupled}.
For a two-species nonlocal PDE with symmetric interactions and an associated
energy principle, see \citet{giunta2024weaklynonlinear}.

\paragraph{Product family and product metric.}
Let $\Omega_i\subset\mathbb R^{d_i}$ and $\Omega=\prod_{i=1}^m\Omega_i$.  We consider the
product family
\begin{equation}
    \mathcal M
    \coloneqq
    \left\{
        \mu=\bigotimes_{i=1}^m\mu_i
        \;\middle|\;
        \mu_i\in\mathcal P_2(\Omega_i)
    \right\}.
\end{equation}
For $\mu=\bigotimes_i\mu_i$ and $\nu=\bigotimes_i\nu_i$, we equip $\mathcal M$ with the
product Wasserstein metric
\begin{equation}
    \label{eq:product-Wasserstein-metric}
    W_{2,\mathcal M}(\mu,\nu)^2
    \coloneqq
    \sum_{i=1}^m W_2(\mu_i,\nu_i)^2 .
\end{equation}
An admissible tangent vector on $\mathcal M$ is a coordinate-wise vector field
$v=(v_1,\ldots,v_m)$ with $v_i\in L^2(\mu_i;\mathbb R^{d_i})$, identified with
$v(x)=(v_1(x_1),\ldots,v_m(x_m))$ for $x=(x_1,\ldots,x_m)$.  The product inner product is
\begin{equation}
    \langle u,v\rangle_{\mu,\mathcal M}
    \coloneqq
    \sum_{i=1}^m \int_{\Omega_i} \langle u_i(x_i),v_i(x_i)\rangle\,\mu_i(\dd x_i)
    =
    \int_{\Omega}\sum_{i=1}^m \langle u_i(x_i),v_i(x_i)\rangle\,\mu(\dd x).
\end{equation}
The corresponding product perturbation keeps the measure in $\mathcal M$:
\begin{equation}
    (\Id+h v)_\sharp\mu
    =
    \bigotimes_{i=1}^m (\Id_i+h v_i)_\sharp\mu_i .
\end{equation}

\paragraph{Restricted Wasserstein gradient.}
Let $\overline F(\mu_1,\ldots,\mu_m)\coloneqq F(\bigotimes_i\mu_i)$.  For
$\mu=\bigotimes_i\mu_i$, write $\mu_{-i}=\bigotimes_{j\neq i}\mu_j$ and
$x_{-i}=(x_j)_{j\neq i}$.  If $F$ admits a first variation, then the partial first
variation of $\overline F$ with respect to $\mu_i$ is, up to an additive constant,
\begin{equation}
    \frac{\delta \overline F}{\delta \mu_i}(\mu_1,\ldots,\mu_m;x_i)
    =
    \int_{\Omega_{-i}}
    \frac{\delta F}{\delta\mu}(\mu;x_i,x_{-i})\,\mu_{-i}(\dd x_{-i}).
\end{equation}
Consequently, the product Wasserstein gradient is
\begin{equation}
    g_i^\mu(x_i)
    \coloneqq
    \nabla_{x_i}\frac{\delta \overline F}{\delta \mu_i}(\mu_1,\ldots,\mu_m;x_i)
    =
    \int_{\Omega_{-i}}
    \left(\nabla_\mu F(\mu;x_i,x_{-i})\right)_i\,\mu_{-i}(\dd x_{-i}),
    \qquad i=1,\ldots,m,
\end{equation}
whenever differentiation and integration can be interchanged.  Equivalently,
$g^\mu=(g_1^\mu,\ldots,g_m^\mu)$ is the orthogonal projection of the full Wasserstein
gradient $\nabla_\mu F(\mu)$ onto coordinate-wise product tangent directions.
Indeed, for every admissible $v=(v_1,\ldots,v_m)$,
\begin{equation}
    \frac{\dd}{\dd h}\Big|_{h=0}F\bigl((\Id+h v)_\sharp\mu\bigr)
    =
    \sum_{i=1}^m \int_{\Omega_i}\langle g_i^\mu(x_i),v_i(x_i)\rangle\,\mu_i(\dd x_i).
\end{equation}

\paragraph{Restricted WGF.}
The Wasserstein gradient flow of $F$ restricted to $\mathcal M$ is the product curve
$\mu_t=\bigotimes_i\mu_{i,t}$ satisfying
\begin{equation}
    \partial_t\mu_{i,t}
    + \nabla_{x_i}\cdot\left(v_{i,t}\mu_{i,t}\right)=0,
    \qquad
    v_{i,t}(x_i)=-g_i^{\mu_t}(x_i),
    \qquad i=1,\ldots,m .
\end{equation}
Equivalently, as a continuity equation on $\Omega$, this reads
\begin{equation}
    \partial_t\mu_t
    + \sum_{i=1}^m \nabla_{x_i}\cdot\left(v_{i,t}(x_i)\mu_t\right)=0.
\end{equation}
Formally, the characteristic chain rule gives the dissipation identity
\begin{equation}
    \frac{\dd}{\dd t}F(\mu_t)
    =
    -\sum_{i=1}^m \int_{\Omega_i}\|g_i^{\mu_t}(x_i)\|^2\,\mu_{i,t}(\dd x_i)
    \le 0.
\end{equation}
In the local matrix-factorization result where this identity is used, its
regularity premises are imposed on the characteristic solution and the
identity is justified for almost every time in
\cref{prop:reduced-dissipation-antisymmetric-branch,eq:mf-local-energy-dissipation}.
